\documentclass[12pt]{article}
\usepackage[utf8]{inputenc}
\usepackage[T1]{fontenc}
\usepackage{hyperref}
\title{Warungu Language Resources for K-12 Curriculum Development}
\author{Compiled by Claude Sonnet 5.5 (claude-sonnet-5-5) for Andr\'as Kornai}
\date{}
\begin{document}
\maketitle

\section{Introduction}
Warungu (ISO 639-3 \texttt{wrg}, Glottocode waru1264) is a dormant Pama-Nyungan language of North Queensland, Australia. Searches show that sources use several spellings (Warungu, Warrongo, Waruŋu), and that search results mix them with Gugu-Badhun and with the unrelated Warumungu of the Northern Territory. Care is needed to match each source to the right language. Wikipedia says the last native speaker, Alf Palmer, died in 1981. It also reports a revival effort, including a first language workshop run by Goondaloo Aboriginal Corporation in 2024. No school teaching was found.

\section{Primary Educational Resources}
No primers, textbooks, or school programmes were found in the search results.

\section{Linguistic and Reference Materials}
\begin{itemize}
\item Tasaku Tsunoda, 1974, \emph{A Grammar of the Waruŋu language, North Queensland} (Monash University MA thesis), listed in Glottolog: \url{https://glottolog.org/resource/reference/id/24813}
\item Tsunoda's later book \emph{A grammar of Warrongo} (De Gruyter Mouton, 2011) is listed in a library record: \url{https://bvbvf1.bib-bvb.de/vufind/Record/DE-604.BV039822751/Details}
\item Wikipedia article on the Warrongo language: \url{https://en.wikipedia.org/wiki/Warrongo_language}
\item Wikipedia article on the Warrongo people: \url{https://en.wikipedia.org/wiki/Warrongo_people}
\item Endangered Languages Project entry: \url{https://endangeredlanguages.com/node/94963}
\item ASJP wordlist entry: \url{https://asjp.clld.org/languages/WARUNGU}
\end{itemize}
Peter Sutton's 1973 thesis \emph{Gugu-Badhun and its Neighbours: a Linguistic Salvage Study} was also mentioned as relevant to neighbouring varieties. No URL for it was seen in the search results.

\section{Religious and Community Resources}
The Goondaloo Aboriginal Corporation was reported (via Wikipedia) to have held a language workshop in 2024. No website or contact details were found.

\section{Digital and Audio Resources}
None found.

\section{Curriculum Development Context}
The language is documented mainly through academic grammars and a short word list. No classroom materials exist in what was found. A heritage curriculum would have to build from Tsunoda's grammar and any archival material, in partnership with the community and with the Goondaloo revival work. Closely related neighbouring languages may offer supplementary material, but they should be treated as separate.
\end{document}